\chapter{Formats de fichiers}
\label{chap:formats}

AirfoilTools sait lire et écrire plusieurs formats de profil, qui
se répartissent en deux familles : les formats \emph{points
discrets} (issus du milieu aéronautique) et les formats
\emph{splines} propres à AirfoilTools.

\section{Formats points discrets}

\subsection{Format Selig (\chemin{.dat})}

Format historique le plus répandu, popularisé par la base de
données de l'Université de l'Illinois (UIUC). C'est le format
par défaut sur \chemin{airfoiltools.com}.

\begin{itemize}
    \item Première ligne : nom du profil (libre).
    \item Lignes suivantes : couples \texttt{x y} séparés par
        des espaces.
    \item Parcours en \emph{convention Selig} : du bord de fuite
        ($x=1$) vers l'extrados (x décroissant), au bord
        d'attaque ($x=0$), puis vers l'intrados (x croissant)
        et retour au bord de fuite.
\end{itemize}

\begin{lstlisting}[caption={Exemple de fichier Selig (NACA~2412)}]
NACA 2412
   1.000000   0.001260
   0.985700   0.003930
   0.971440   0.006580
   ...
   0.014300   0.018780
   0.000000   0.000000
   0.014300  -0.014050
   ...
   0.985700  -0.001280
   1.000000  -0.001260
\end{lstlisting}

\subsection{Format Lednicer (\chemin{.dat})}

Variante du Selig où extrados et intrados sont stockés
séparément, chacun parcouru du bord d'attaque vers le bord de
fuite. Une ligne d'en-tête indique le nombre de points de chaque
côté.

\begin{lstlisting}[caption={Exemple de fichier Lednicer}]
NACA 2412
65.0   65.0

   0.000000   0.000000
   0.001425   0.005700
   ...
   1.000000   0.001260

   0.000000   0.000000
   0.001425  -0.005720
   ...
   1.000000  -0.001260
\end{lstlisting}

\subsection{Format CSV (\chemin{.csv})}

Format tabulaire avec en-tête \chemin{x,y}. Pratique pour
l'import dans des feuilles de calcul ou pour le partage avec des
outils non-aéronautiques.

\begin{lstlisting}[caption={Exemple de fichier CSV}]
x,y
1.000000,0.001260
0.985700,0.003930
...
\end{lstlisting}

\section{Formats splines AirfoilTools}

\subsection{Format \chemin{.bspl}}

Sauvegarde complète d'un profil en mode \textbf{spline
multi-segment}. Préserve l'intégralité de la définition Bézier :
points de contrôle, degrés, continuités aux jonctions, paramètres
d'échantillonnage.

C'est le format \textbf{recommandé} pour conserver le travail
d'édition d'un profil et le ré-ouvrir sans perte d'information.

\begin{noteinfo}
Le format \chemin{.bspl} est un format texte structuré (JSON
ou similaire) propre à AirfoilTools. Il n'est pas
interopérable avec d'autres logiciels de CAO ou
d'aérodynamique. Pour une diffusion externe, ré-exportez en
\chemin{.dat} (Selig) après avoir ajusté l'échantillonnage.
\end{noteinfo}

\subsection{Format \chemin{.bez}}

Variante simplifiée du \chemin{.bspl} pour les profils définis
par une seule courbe de Bézier par côté (mode mono-segment).
Conservé pour compatibilité avec les versions antérieures.

\section{Génération NACA}

L'application sait générer à la volée des profils NACA à partir
de leur désignation, sans passer par un fichier :

\begin{itemize}
    \item NACA 4 chiffres (ex. \texttt{2412}, \texttt{0012}) :
        cambrure maximale (\%), position de la cambrure max
        (dixièmes), épaisseur relative (\%).
    \item NACA 5 chiffres (ex. \texttt{23012}) : profil cambré
        avec ligne moyenne réflexe.
\end{itemize}

Cette génération se fait automatiquement au démarrage de
l'application (NACA~2412 et NACA~0012 par défaut). Pour générer
d'autres NACA en cours de session, utilisez l'API Python ou
préchargez un fichier \chemin{.dat} obtenu ailleurs.

\section{Tableau récapitulatif}

\begin{tabularx}{\linewidth}{l c c X}
    \toprule
    \textbf{Format} & \textbf{Lecture} & \textbf{Écriture}
        & \textbf{Préserve la spline ?} \\
    \midrule
    Selig (\chemin{.dat})    & oui & oui & non \\
    Lednicer (\chemin{.dat}) & oui & oui & non \\
    CSV (\chemin{.csv})      & oui & oui & non \\
    Spline (\chemin{.bspl})  & oui & oui & \textbf{oui} \\
    Bézier (\chemin{.bez})   & oui & oui & \textbf{oui} (mono-segment) \\
    \bottomrule
\end{tabularx}

\begin{important}
Lorsque vous sauvegardez un profil au format \chemin{.dat} ou
\chemin{.csv}, \textbf{l'information de spline est perdue}~:
seule la liste des points discrets échantillonnés est écrite.
Pour préserver intégralement votre travail, sauvegardez
toujours en \chemin{.bspl}.
\end{important}
